\section{Introduction}\label{sec:intro}

Unsupervised classification partitions a data set into groups without labelled
examples, and in doing so it must answer a prior question: how many groups are
there? The choice of the number of clusters $K$ governs every downstream
interpretation, and in the analysis of remotely sensed imagery, where $K$ is
the number of land-cover classes an analyst will map, it is made routinely
and consequentially. A mature toolkit exists to make it: the gap statistic
\citep{Tibshirani2001}, the average silhouette \citep{Rousseeuw1987}, the
Calinski--Harabasz \citep{CalinskiHarabasz1974} and Davies--Bouldin
\citep{DaviesBouldin1979} indices each score a partition by its internal
geometry and select the $K$ that optimises the score.
\revtext{The stakes are concrete. An analyst mapping land cover from a
satellite scene must fix $K$ before any unsupervised classifier runs, and
every downstream product (class areas, change statistics, sampling strata,
training masks) inherits the choice. Set $K$ too high and single crop fields
fragment into spurious sub-classes that propagate into those products; set it
too low and real classes merge. The same prior decision arises when
delineating geochemical domains, zoning a mineral deposit, or partitioning a
scene into regions for stratified field sampling: in each case the analyst
must first ask how many classes the data actually support.}

\revtext{These criteria share a blind spot: none of them takes the dependence
between observations into account.} Their reference levels, the uniform null
of the gap statistic and
the implicit baseline of a silhouette, describe what the geometry would look
like were the samples exchangeable draws. Spatial data are not exchangeable.
Neighbouring pixels are autocorrelated \citep{Tobler1970}, so the effective
number of independent observations is far below the nominal pixel count
\citep{Cressie1993}, and a smooth illumination or topographic gradient spreads
the feature cloud along a manifold that an independence null cannot
distinguish from genuine groups. The consequence, which we quantify in
Sections~\ref{sec:simulation}--\ref{sec:application}, is systematic
\emph{over-detection}\revtext{ (throughout, the partitions the criteria
score are produced by $k$-means, the dominant choice in this setting;
statements about a criterion's behaviour are therefore statements about the
criterion--$k$-means pair, a dependence Section~\ref{subsec:methods} makes
explicit)}. On synthetic fields containing no discrete classes the
gap statistic reports on average between seven and eight; on real imagery
the silhouette, Calinski--Harabasz and Davies--Bouldin indices fragment a
single, homogeneous crop field into several sub-classes in every instance we
examined. A criterion that cannot recognise the absence of clusters is a
precarious foundation for deciding their number.

Correcting this is not a matter of swapping in a better null. Because the
within-cluster dispersion that these criteria minimise is a variance-based
quantity, a reference matched to the data's spatial covariance structure,
its variogram, reproduces the same dispersion and explains the clusters
away. Distinguishing real classes from autocorrelation therefore requires
information beyond the covariance, namely the \emph{multimodality} of the feature
distribution; and distinguishing real classes from a smooth trend requires
separating a continuous drift from discrete structure, a separation the
geostatistical decomposition into trend and residual makes natural. The idea of
calibrating a clustering criterion against an autocorrelation-aware null is
itself not new. \citet{HennigLin2015} do exactly this for discrete
presence--absence data, but it has not been carried to the continuous-field
setting of gridded imagery, nor has the confounding of trend with class been
addressed in selecting $K$.

This paper supplies both. We recast the question \emph{how many classes?} as a
sequence of homogeneity decisions: at each node we test whether a region is a
single class or instead carries \emph{multimodality in excess of spatial
autocorrelation}, and make that test valid on autocorrelated fields by
calibrating it at the effective sample size. The
construction rests on three geostatistical devices. The first is the effective
sample size itself \citep{Cressie1993,Dutilleul1993}, which sets the
calibration; it is evaluated at the indicator integral range appropriate to a
functional of the distribution, not at the field integral range that governs
the mean (Section~\ref{subsec:sgs}). The second is a \emph{local} estimate of
the within-class autocorrelation range, in the manner of declustering
\citep{IsaaksSrivastava1989,Deutsch1998}, which prevents the between-class
mosaic from corrupting that calibration. The third is an adaptive
trend--residual decomposition, which removes a smooth drift only when it is
statistically resolvable. A recursive partitioning driven by the calibrated dip statistic
\citep{HartiganHartigan1985} returns $K$. Because the test is tuned for
specificity, its negatives are trustworthy: when it declares a region
homogeneous that verdict can be relied on, which makes ESS-Dip as useful as a
homogeneity guard or a stopping rule for region-growing as it is as a
conservative estimator of $K$. Our contributions are:

\begin{itemize}
  \item We document, on controlled ground truth, that the standard
    number-of-clusters criteria over-detect under spatial autocorrelation
    without exception, failing to recognise class-free scenes in every one of
    $75$ realisations (Section~\ref{sec:simulation}).
  \item We propose \emph{ESS-Dip}, an effective-sample-size-calibrated test of
    cluster homogeneity (whether a region holds discrete structure beyond its
    spatial autocorrelation), combining a within-class range from local
    range estimation with adaptive trend removal (Section~\ref{sec:method}).
    Applied recursively it returns a conservative $K$, and we state an
    asymptotic false-split-control property (Proposition~\ref{prop:fpr}) that
    the method confirms empirically. A recall-oriented variant, ESS-Dip-R,
    trades the distribution-free uniform null for the model's Gaussian null to
    recover more structure at the same specificity.
  \item We evaluate \revtext{sixteen} criteria on a factorial simulation with known $K$\revtext{,
    including the most direct spatial adaptation of a standard index (the gap
    statistic with an autocorrelation-matched reference)},
    including the calibrated-validity-index method of \citet{HennigLin2015} that
    ESS-Dip generalises; ESS-Dip reaches unit specificity, its recall variant
    ESS-Dip-R attains the best balanced score at the same specificity, and an
    ablation attributes the gain to the local range and the adaptive detrending
    (Section~\ref{sec:simulation}).
  \item We confirm the behaviour on three real scenes spanning hyperspectral
    and multispectral sensing (Indian Pines, Salinas, and a Sentinel-2 scene
    labelled by ESA WorldCover), where ESS-Dip identifies single land-cover
    regions as homogeneous where the standard criteria do not
    (Section~\ref{sec:application}).
\end{itemize}

The remainder of the paper is organised as follows.
Section~\ref{sec:background} reviews number-of-clusters criteria, the
geostatistical notions of effective sample size and declustering, and
modality-based clustering, and positions the contribution.
Section~\ref{sec:method} develops the method and its calibration.
Sections~\ref{sec:simulation} and~\ref{sec:application} report the simulation
study and the application to land-cover imagery. Section~\ref{sec:discussion}
discusses the precision--recall trade-off and limitations, and
Section~\ref{sec:conclusions} concludes. Code and data are released to
reproduce every result.
